\chapter{Array in Java}
\label{chap:array-java}

Dopo le strutture dati e i riferimenti introdotti nei
\cref{chap:strutture-dati-adt,chap:introduzione-java}, si esamina
l'uso degli array in Java: accesso agli elementi, elaborazione tramite
metodi, array bidimensionali e organizzazione di dati correlati.\footnote{Materiale
di riferimento: S. Canazza, \emph{Java: Arrays},
\nolinkurl{2026_FI_12_canazza_Java_Array.key.pdf}.}

\section{Struttura, costruzione e accesso agli array}
\label{sec:fi-java-array-fondamenti}

\begin{definition}[Array in Java]
\label{def:fi-java-array}
Un \emph{array} è un oggetto che contiene una sequenza di elementi
omogenei, individuati mediante indici interi consecutivi.
La sua lunghezza è stabilita alla costruzione e non può cambiare.
\end{definition}

Il concetto astratto e le proprietà dell'accesso diretto sono già
stati esaminati nella \cref{sec:fi-strutture-array}.
In Java un array può contenere valori primitivi oppure riferimenti
a oggetti; gli oggetti indicati da tali riferimenti sono distinti
dall'array che li contiene.

\subsection{Dichiarazione, costruzione e indici}
\label{sec:fi-java-array-costruzione}

Si dichiara una variabile di tipo \texttt{double[]} e si costruisce
un array di dieci elementi mediante \texttt{new}:

\begin{unipdcode}[language=Java]
double[] valori = new double[10];
valori[3] = 2.5;
double elemento = valori[3];
\end{unipdcode}

\texttt{valori} contiene un riferimento all'array;
\texttt{valori[3]} ne individua il quarto elemento. Gli indici
ammessi appartengono all'intervallo
\(0\le i<\texttt{valori.length}\). Un indice negativo o non inferiore
alla lunghezza provoca \texttt{ArrayIndexOutOfBoundsException},
un'eccezione non controllata
(\cref{sec:fi-java-file-checked}).
La lunghezza richiesta a \texttt{new} non può essere negativa:
in tal caso viene lanciata \texttt{NegativeArraySizeException}.

\subsection{Lunghezza e inizializzazione}
\label{sec:fi-java-array-inizializzazione}

\texttt{length} è un campo \texttt{public final}, non un metodo:
\texttt{valori.length} restituisce il numero di elementi,
mentre assegnargli un nuovo valore è un errore di compilazione.
La lunghezza può essere zero. Gli elementi di un array costruito
con \texttt{new} assumono i valori iniziali previsti dal tipo:
zero per i tipi numerici, \texttt{false} per \texttt{boolean},
\texttt{\textbackslash u0000} per \texttt{char} e \texttt{null} per i riferimenti.

\begin{unipdcode}[language=Java]
int[] a = new int[3];
a[0] = 2; a[1] = 3; a[2] = 5;
int[] b = {2, 3, 5};
int[] c = new int[] {2, 3, 5};
\end{unipdcode}

Le tre costruzioni producono array distinti con lo stesso contenuto.
La forma abbreviata con le sole graffe è ammessa nella dichiarazione
con inizializzazione, non in un assegnamento ordinario successivo.
Per un array di riferimenti, l'inizializzazione delle celle a
\texttt{null} non costruisce automaticamente gli oggetti indicati.

\subsection{Parametri, valori di ritorno e copia}
\label{sec:fi-java-array-parametri}

Un parametro di tipo \texttt{int[]} riceve una copia del riferimento,
secondo il passaggio per valore della
\cref{def:fi-java-file-passaggio-valore}. Il metodo può modificare
le celle dell'array ricevuto, ma riassegnare il parametro non cambia
la variabile del chiamante. Un metodo può anche restituire un
riferimento a un array:

\begin{unipdcode}[language=Java]
public static int somma(int[] v) {
    if (v == null) throw new IllegalArgumentException();
    int totale = 0;
    for (int i = 0; i < v.length; i++) totale += v[i];
    return totale;
}

public static int[] primiTre() {
    return new int[] {2, 3, 5};
}
\end{unipdcode}

L'assegnamento \texttt{int[] y = x;} copia il riferimento, non
l'array: modificare \texttt{y[0]} modifica anche \texttt{x[0]}.
Per ottenere un array indipendente occorre copiarne gli elementi.
Si possono usare un ciclo, \texttt{System.arraycopy()} oppure
\texttt{clone()}:

\begin{unipdcode}[language=Java]
int[] x = {2, 3, 5};
int[] y = x;                       // stesso array
int[] z = x.clone();               // nuovo array
int[] w = new int[x.length];
System.arraycopy(x, 0, w, 0, x.length);
\end{unipdcode}

I cinque parametri di \texttt{System.arraycopy()} identificano
array sorgente e indice iniziale, array destinazione e indice
iniziale, e numero di elementi da copiare. Gli
intervalli devono essere validi e i tipi compatibili; il metodo
gestisce anche intervalli sovrapposti nello stesso array.
Il metodo \texttt{clone()} degli array restituisce un riferimento
del tipo array appropriato, senza richiedere un cast esplicito.
Entrambe le operazioni effettuano una \emph{copia superficiale}:
per un array di oggetti, copiano i riferimenti, non gli oggetti.

\begin{figure}[htbp]
  \centering
  \begin{tikzpicture}[unipd/diagram,x=1cm,y=1cm]
  \node[unipd/node,minimum width=1.5cm]
    (x) at (-3.4,0.45) {\texttt{x}};
  \node[unipd/node,minimum width=1.5cm]
    (y) at (-3.4,-0.45) {\texttt{y}};

  \node[unipd/node,minimum width=3cm,
    minimum height=1.5cm,align=center]
    (a) at (0,0) {Array originale\\\texttt{[2, 3, 5]}};

  \node[unipd/node,minimum width=1.5cm]
    (z) at (3.4,-2.65) {\texttt{z}};
  \node[unipd/node,minimum width=3cm,align=center]
    (b) at (0,-2.65) {Copia indipendente\\\texttt{[2, 3, 5]}};

  \draw[unipd/directed edge]
    (x.east) -- ([yshift=0.45cm]a.west);
  \draw[unipd/directed edge]
    (y.east) -- ([yshift=-0.45cm]a.west);
  \draw[unipd/directed edge]
    (z.west) -- (b.east);

  \draw[unipd/directed edge,densely dashed]
    (a.south) -- node[midway,right=6pt,
      font=\scriptsize,fill=white,inner sep=2pt]
      {copia degli elementi} (b.north);
\end{tikzpicture}
  \caption{Due variabili possono riferirsi allo stesso array;
  la copia degli elementi produce un oggetto distinto.}
  \label{fig:fi-java-array-copie}
\end{figure}

\clearpage

\section{Il ciclo \texttt{for-each}}
\label{sec:fi-java-array-for-each}

Il \emph{for-each}, o \emph{enhanced for}, visita gli elementi
di un array in ordine senza utilizzare un indice esplicito:

\begin{unipdcode}[language=Java]
int[] voti = {24, 27, 30};
int somma = 0;
for (int voto : voti) {
    somma += voto;
}
System.out.println(somma);  // 81
\end{unipdcode}

Ad ogni iterazione, \texttt{voto} riceve il valore dell'elemento
corrente. Assegnare un nuovo valore a \texttt{voto} non modifica
la cella originaria. Per accedere agli indici, sostituire gli
elementi o visitare soltanto determinate posizioni, è preferibile
il \texttt{for} tradizionale
(\cref{sec:fi-java-cicli-for}).
Il \texttt{for-each} è utilizzabile anche sulle collezioni che
implementano \texttt{Iterable}; gli array non devono implementare
tale interfaccia per essere visitati in questa forma.

\section{Ridimensionamento mediante copia}
\label{sec:fi-java-array-resize}

La lunghezza di un array non è modificabile. Per ottenere una
lunghezza diversa si costruisce un nuovo array, si trasferiscono
i dati da conservare e si restituisce il nuovo riferimento.
Il metodo statico seguente è usato anche nella
\cref{sec:fi-java-array-algoritmi}:

\begin{unipdcode}[language=Java]
public static int[] resize(int[] original, int nuovaLunghezza) {
    if (original == null || nuovaLunghezza < 0)
        throw new IllegalArgumentException();
    int[] copia = new int[nuovaLunghezza];
    System.arraycopy(original, 0, copia, 0,
                     Math.min(original.length, nuovaLunghezza));
    return copia;
}
\end{unipdcode}


Il numero di elementi conservati è
\(\min(\texttt{original.length},\texttt{nuovaLunghezza})\).
Se il nuovo array è più lungo, le celle aggiuntive mantengono
l'inizializzazione predefinita; se è più corto, gli elementi
finali vengono scartati. Il chiamante deve usare il riferimento restituito; gli altri
eventuali riferimenti all'array originale continuano a indicare
l'oggetto precedente:

\begin{unipdcode}[language=Java]
int[] valori = {1, 3, 7};
valori = ArrayAlgs.resize(valori, 5);  // [1, 3, 7, 0, 0]
\end{unipdcode}

\clearpage

\section{Algoritmi su array: la classe \texttt{ArrayAlgs}}
\label{sec:fi-java-array-algoritmi}

La classe di utilità \texttt{ArrayAlgs} raccoglie metodi
\texttt{static} per elaborare array di interi
(\cref{sec:fi-java-cicli-metodi-statici}). Ogni listato seguente
mostra un singolo metodo della stessa classe; non ne viene
ripetuta la dichiarazione. Le operazioni che cambiano la lunghezza
restituiscono un nuovo array.

\subsection{Generazione di numeri pseudocasuali}
\label{sec:fi-java-array-casuale}

\texttt{Math.random()} restituisce un valore \texttt{double}
pseudocasuale nell'intervallo \([0,1)\). Il metodo costruisce un
array di lunghezza \texttt{length} con interi nell'intervallo
\emph{inclusivo} \([\texttt{min},\texttt{max}]\):

\begin{unipdcode}[language=Java]
public static int[] randomIntArray(int length, int min, int max) {
    if (length < 0 || min > max)
        throw new IllegalArgumentException();
    int[] result = new int[length];
    long intervallo = (long) max - min + 1;
    for (int i = 0; i < result.length; i++)
        result[i] = (int) (min + (long) (Math.random() * intervallo));
    return result;
}
\end{unipdcode}


La differenza viene calcolata con \texttt{long} per evitare
l'overflow nell'estensione dell'intervallo; i valori restituiti
restano di tipo \texttt{int}. Lunghezze negative e intervalli
invertiti non sono ammessi.

\subsection{Rappresentazione testuale}
\label{sec:fi-java-array-stampa}

Stampare direttamente un riferimento a un array con
\texttt{System.out.println(v)} non ne elenca gli elementi.
Il metodo seguente costruisce una rappresentazione testuale:

\begin{unipdcode}[language=Java]
public static String printArray(int[] v) {
    if (v == null) throw new IllegalArgumentException();
    String testo = "[";
    for (int i = 0; i < v.length; i++) {
        if (i > 0) testo += ", ";
        testo += v[i];
    }
    return testo + "]";
}
\end{unipdcode}


\begin{unipdcode}[language=Java]
int[] v = {2, 3, 5};
System.out.println(ArrayAlgs.printArray(v)); // [2, 3, 5]
\end{unipdcode}

\clearpage

\subsection{Rimozione con e senza conservazione dell'ordine}
\label{sec:fi-java-array-remove}

La rimozione restituisce un array con un elemento in meno. Se l'ordine non è rilevante, l'elemento rimosso
può essere sostituito dall'ultimo; altrimenti gli elementi
successivi devono conservare l'ordine relativo.

\noindent
\begin{minipage}[t]{0.48\linewidth}
\vspace{0pt}
\textbf{Senza conservare l'ordine}

\begin{unipdcode}[language=Java,
  basicstyle=\ttfamily\scriptsize, numbers=none,
  xleftmargin=0.6em, xrightmargin=0.6em,
  framexleftmargin=0.2em]
public static int[] remove(
        int[] v, int index) {
    if (v == null || index < 0
            || index >= v.length)
        throw new IllegalArgumentException();
    int[] out = resize(v, v.length - 1);
    if (index < out.length)
        out[index] = v[v.length - 1];
    return out;
}
\end{unipdcode}

\end{minipage}%
\hfill
\begin{minipage}[t]{0.48\linewidth}
\vspace{0pt}
\textbf{Conservando l'ordine}

\begin{unipdcode}[language=Java,
  basicstyle=\ttfamily\scriptsize, numbers=none,
  xleftmargin=0.6em, xrightmargin=0.6em,
  framexleftmargin=0.2em]
public static int[] removeSorted(
        int[] v, int index) {
    if (v == null || index < 0
            || index >= v.length)
        throw new IllegalArgumentException();
    int[] out = new int[v.length - 1];
    for (int i = 0; i < index; i++)
        out[i] = v[i];
    for (int i = index; i < out.length; i++)
        out[i] = v[i + 1];
    return out;
}
\end{unipdcode}

\end{minipage}

\begin{unipdcode}[language=Java]
int[] v = {1, 2, 3, 4, 5};
int[] a = ArrayAlgs.remove(v, 1);        // [1, 5, 3, 4]
int[] b = ArrayAlgs.removeSorted(v, 1);  // [1, 3, 4, 5]
\end{unipdcode}

\texttt{removeSorted()} conserva l'ordine originario, ma non
esegue un ordinamento. Entrambi i metodi restituiscono un nuovo
array: la copia richiede \(\Theta(n)\) tempo e spazio per \(n\)
elementi. La differenza riguarda i movimenti necessari dopo la
rimozione: una sostituzione con l'ultimo elemento, oppure lo
scorrimento della parte successiva della sequenza.

\subsection{Inserimento in una posizione}
\label{sec:fi-java-array-insert}

Per inserire un elemento all'indice \texttt{index}, compreso
fra zero e la lunghezza originaria, si crea un array con una
cella aggiuntiva. Gli elementi da \texttt{index} in poi vengono spostati verso
indici maggiori, procedendo dall'ultimo verso il primo per
non sovrascrivere valori ancora necessari:

\begin{unipdcode}[language=Java]
public static int[] insert(int[] v, int index, int value) {
    if (v == null || index < 0 || index > v.length)
        throw new IllegalArgumentException();
    int[] out = resize(v, v.length + 1);
    for (int i = v.length; i > index; i--)
        out[i] = out[i - 1];
    out[index] = value;
    return out;
}
\end{unipdcode}


\begin{unipdcode}[language=Java]
int[] v = {1, 2, 3, 4, 5};
v = ArrayAlgs.insert(v, 2, 7);  // [1, 2, 7, 3, 4, 5]
\end{unipdcode}

Gli elementi vengono trasferiti verso \emph{indici maggiori},
ma sono elaborati a partire dall'ultimo. Il costo temporale
e lo spazio aggiuntivo sono \(\Theta(n)\).

\clearpage

\subsection{Ricerca lineare}
\label{sec:fi-java-array-linear-search}

La ricerca lineare già studiata nella \cref{sec:fi-ricerca-lineare}
esamina gli elementi nell'ordine e restituisce il primo indice
corrispondente, oppure \(-1\) se il valore non è presente:

\begin{unipdcode}[language=Java]
public static int linearSearch(int[] v, int value) {
    if (v == null) throw new IllegalArgumentException();
    for (int i = 0; i < v.length; i++)
        if (v[i] == value) return i;
    return -1;
}
\end{unipdcode}


Nel caso peggiore vengono effettuati \(n\) confronti:
complessità temporale \(\Theta(n)\) e spazio ausiliario
\(\Theta(1)\).

\subsection{Ricerca del massimo e del minimo}
\label{sec:fi-java-array-estremi}

Per un array non vuoto, il primo elemento costituisce il
candidato iniziale. I successivi aggiornano il candidato
quando si trova un valore maggiore o minore:

\noindent
\begin{minipage}[t]{0.48\linewidth}
\vspace{0pt}
\textbf{Massimo}

\begin{unipdcode}[language=Java,
  basicstyle=\ttfamily\scriptsize, numbers=none,
  xleftmargin=0.6em, xrightmargin=0.6em,
  framexleftmargin=0.2em]
public static int findMax(int[] v) {
    if (v == null || v.length == 0)
        throw new IllegalArgumentException();
    int max = v[0];
    for (int i = 1; i < v.length; i++)
        if (v[i] > max) max = v[i];
    return max;
}
\end{unipdcode}

\end{minipage}%
\hfill
\begin{minipage}[t]{0.48\linewidth}
\vspace{0pt}
\textbf{Minimo}

\begin{unipdcode}[language=Java,
  basicstyle=\ttfamily\scriptsize, numbers=none,
  xleftmargin=0.6em, xrightmargin=0.6em,
  framexleftmargin=0.2em]
public static int findMin(int[] v) {
    if (v == null || v.length == 0)
        throw new IllegalArgumentException();
    int min = v[0];
    for (int i = 1; i < v.length; i++)
        if (v[i] < min) min = v[i];
    return min;
}
\end{unipdcode}

\end{minipage}

Il caso vuoto viene rifiutato perché non possiede un massimo
né un minimo. Entrambi i metodi hanno costo \(\Theta(n)\) e
utilizzano \(\Theta(1)\) spazio ausiliario, secondo la
\cref{sec:fi-ricerca-estremi}.

\section{Array bidimensionali}
\label{sec:fi-java-array-bidimensionali}

Un array bidimensionale in Java è un \emph{array di array}.
L'espressione \texttt{v[i][j]} seleziona l'elemento in posizione
\texttt{j} della riga \texttt{i}; la prima lunghezza è
\texttt{v.length}, quella della riga \texttt{i} è
\texttt{v[i].length}.

\begin{unipdcode}[language=Java]
int[][] potenze = new int[4][5];
for (int i = 0; i < potenze.length; i++) {
    for (int j = 0; j < potenze[i].length; j++) {
        potenze[i][j] = (int) Math.pow(i + 1, j + 1);
    }
}
System.out.println(potenze[2][3]);  // 81, ossia 3^4
\end{unipdcode}

\clearpage

La scansione utilizza i cicli annidati della
\cref{sec:fi-java-cicli-for-annidati}. In questo caso gli indici
di riga sono \(0,\ldots,3\) e quelli di colonna \(0,\ldots,4\). L'accesso a una riga o
a una sua cella con indice non valido provoca un'eccezione.
Le righe possono avere lunghezze differenti:

\begin{unipdcode}[language=Java]
int[][] righe = {{1, 2, 3}, {4}, {5, 6}};
System.out.println(righe[0].length); // 3
System.out.println(righe[1].length); // 1
\end{unipdcode}

Non esiste quindi, in generale, un unico numero di colonne.
Prima di usare \texttt{v[0].length} occorre verificare che
l'array esterno non sia vuoto; una riga nulla non può
essere indicizzata. Le scansioni utilizzano pertanto
\texttt{v[i].length}, non la lunghezza della prima riga.

\section{Argomenti della riga di comando}
\label{sec:fi-java-array-argomenti}

Il parametro \texttt{String[] args} di \texttt{main}, introdotto
nella \cref{sec:fi-java-primo-programma}, contiene gli argomenti
passati dopo il nome della classe, senza includerlo.
Gli argomenti sono stringhe, anche quando rappresentano numeri:

\begin{unipdcode}[language=Java]
public class Argomenti {
    public static void main(String[] args) {
        System.out.println("Argomenti: " + args.length);
        for (String arg : args)
            System.out.println(arg);
    }
}
\end{unipdcode}

\begin{unipdterminal}
$ java Argomenti 2 33 Hello
Argomenti: 3
2
33
Hello
\end{unipdterminal}

Per convertire una stringa numerica si può usare
\texttt{Integer.parseInt()}, che può lanciare
\texttt{NumberFormatException}
(\cref{sec:fi-java-file-checked}). Le opzioni possono essere
riconosciute convenzionalmente dal prefisso \texttt{-}, per
esempio mediante \texttt{arg.startsWith("-")}; la sintassi
concreta delle opzioni è stabilita dal programma.

\clearpage

\section{Array paralleli e array di oggetti}
\label{sec:fi-java-array-paralleli}

Si parla di \emph{array paralleli} quando elementi allo stesso
indice, in array diversi, descrivono lo stesso dato logico:

\begin{unipdcode}[language=Java]
String[] cognomi = {"Rossi", "Bianchi"};
double[] scritti = {26, 28};
double[] orali = {30, 27};
double media = (scritti[0] + orali[0]) / 2; // 28.0
\end{unipdcode}

L'associazione dipende dall'indice. Inserimenti, eliminazioni
e riordinamenti devono essere coordinati fra tutti gli array;
i metodi richiedono inoltre più parametri per elaborare i
dati dello stesso studente. Una rappresentazione orientata
agli oggetti riunisce invece le proprietà correlate in
un'unica classe (\cref{sec:fi-java-classi-incapsulamento}):

\begin{unipdcode}[language=Java]
class Studente {
    private final String cognome;
    private final double scritto, orale;

    public Studente(String cognome, double scritto, double orale) {
        this.cognome = cognome;
        this.scritto = scritto;
        this.orale = orale;
    }

    public double media() { return (scritto + orale) / 2; }
}
\end{unipdcode}

\begin{unipdcode}[language=Java]
Studente[] studenti = {
    new Studente("Rossi", 26, 30),
    new Studente("Bianchi", 28, 27)
};
System.out.println(studenti[0].media());  // 28.0
\end{unipdcode}

Un nuovo dato riferito allo studente può essere aggiunto alla
classe senza introdurre un ulteriore array da sincronizzare.
L'array di oggetti mantiene l'ordine degli studenti, mentre
ogni oggetto conserva in modo unitario le proprie informazioni.
